\documentclass{article}

% Language setting
% Replace `english' with e.g. `spanish' to change the document language
\usepackage[czech]{babel}
\usepackage{amsthm,thmtools,xcolor,amsmath,amssymb, mathtools}

% Set page size and margins
% Replace `letterpaper' with `a4paper' for UK/EU standard size
\usepackage[a4paper,top=2cm,bottom=2cm,left=3cm,right=3cm,marginparwidth=1.75cm]{geometry}

% Useful packages
\usepackage{enumerate}
\usepackage{graphicx}
\usepackage{tabularx}
\usepackage[colorlinks=true, allcolors=blue]{hyperref}
\graphicspath{images/}

%\usepackage{tikzit}
%\input{default.tikzstyles}

\title{DÚ Lineární algebra -- Sada 8}
\author{Jan Romanovský}

\begin{document}
\maketitle

\textbf{(8.1)} Invariantní prostory budeme hledat jako lineární obal spojení Jordanových řetízků. Spočítáme vlastní čísla a vlastní vektory matice A
$$\lambda_{1,2} = 2, \lambda_3 = -1;\, M_{2} = \text{span}\left\{\left(\begin{array}{c}
  -2 \\ 3 \\ 5
\end{array}\right)\right\}, M_{-1} = \text{span}\left\{\left(\begin{array}{c}
  -1 \\ 0 \\ 1
\end{array}\right)\right\}.$$
Vidíme, že matice má Jordanův tvar, doplníme tedy řetízky, ozn. $J_2, J_{-1}$, pro což nám stačí najít třetí a poslední vektor, který bude prvkem $J_2$, jako řešení SLR
$$\left(\begin{array}{c c c | c}
  -3 &-2& 0& -2\\ -1& 1& -1& 3\\ 2& 3& -1& 5
\end{array}\right);\, J_{2} = \left\{\left(\begin{matrix}
  -2 \\ 3 \\ 5
\end{matrix}\right), \left(
\begin{array}{c}
 -\frac{4}{5} \\
 \frac{11}{5} \\
 0 \\
\end{array}
\right)\right\}, J_{-1}\left\{ \left(\begin{matrix}
  -1 \\ 0 \\ 1
\end{matrix}\right)\right\}.$$

Jediné spojení Jordanových řetízků, které má dva vektory (tedy takové, aby vzniklý podprostor --- lineární obal byl dimenze 2), je samotný řetízek $J_2$, matice zúženého operátor $f_A$ na lineární obal tohoto řetízku bude pak Jordanova buňka příslušná tomuto řetízku

$$[f_A|_{\text{span } J_2}]^{J_2}_{J_2} = \begin{pmatrix}
  2 & 1 \\ 0 & 2
\end{pmatrix}.$$
\end{document}
